%%%%%%%%%%  scan idx 21  |  folio 8  %%%%%%%%%%
qui n'est d'ailleurs autre que
\[
\cl(L_\bullet)=\sum_i(-1)^i\cl^\bullet(L'_i).
\]
où $L'_\bullet$ est comme ci-dessus, de sorte que les $L'_i$ sont bien des Modules localement libres sur $Y$ dont on peut prendre les classes dans $K^\bullet(Y)$.

Ceci posé, on définit un homomorphisme (3.1) chaque fois que $f:X\to Y$ est un morphisme \underline{propre de tor-dimension finie} de schémas noethériens tels que $Y$ admette un Module inversible ample, en posant simplement
\begin{equation}\tag{3.2}
f_*(\cl^\bullet(\underline F))=\cl^\bullet(Rf_*(\underline F))
\end{equation}
pour tout Module localement libre $\underline F$ sur $X$ (de sorte que $Rf_*(\underline F)$ est bien un complexe parfait sur $Y$, et le deuxième membre est défini). Lorsque $X$ également admet un Module inversible ample, alors $K^\bullet(X)$ peut s'interpréter également comme $K(\Parf(X))$, et l'homomorphisme (3.1) n'est autre que celui qui est associé, par fonctorialité, au foncteur exact de catégories triangulées induit par
\[
Rf_*:\Parf(X)\longrightarrow \Parf(Y).
\]

\subsubsection*{3.2.}
Pour donner un sens à (1.2), il reste encore à faire sur $f$ une hypothèse permettant de donner un sens à l'homomorphisme (2.5)
\begin{equation}\tag{3.3}
f_*:\Gr^\bullet(X)_{\mathbb Q}\longrightarrow \Gr^\bullet(Y)_{\mathbb Q},
\end{equation}
i.e.\ assurant que l'homomorphisme (3.1) qu'on vient de définir est (modulo un décalage et modulo torsion) compatible aux filtrations de ces anneaux~; et enfin à définir encore un élément «~fibré tangent relatif virtuel~»
\[
T_f\in K^\bullet(X).
\]
Pour ce dernier, reprenant la définition proposée dans 2.1, on voit qu'il y a lieu de faire une hypothèse assurant que (avec les notations de loc. cit.) le Module conormal $\underline N_{X/X'}$ soit localement libre. L'hypothèse la plus naturelle à cet effet est celle que $i$ soit une «~immersion régulière~», i.e.\ identifie $X$ à un sous-schémas de $X'$ défini par un Idéal qui localement peut être défini par une suite $\underline O_{X'}$-régulière de sections de $\underline O_{X'}$. Cette hypothèse, de nature purement locale, ne dépend pas en fait du choix de la factorisation de $f$ en une immersion suivie d'un morphisme lisse (Exp.\ VIII), et comme plus haut, il en est de même pour la classe $T_f$ défini par la formule (2.1). Lorsque l'hypothèse qu'on vient d'envisager est vérifiée, on dit que le morphisme $f$ est un \underline{morphisme} «~\underline{localement d'insersection complète}~»\footnote{On se permettra également de supprimer le vocable «~localement~» dans cette terminologie, ce qui ne risque d'entrainer aucune confusion.}. Un tel morphisme est d'ailleurs localement de tor-dimension finie.

